\documentclass[12pt]{article}
\usepackage{amsmath, amssymb}
\usepackage{geometry}
\usepackage{setspace}
\geometry{margin=1in}
\onehalfspacing

\title{Lecture 15--17 Scribe \\ Joint Random Variables}
\author{}
\date{}

\begin{document}

\maketitle

\section{Why Study Joint Random Variables}

In real life we often study two random variables at the same time.

Examples:
\begin{itemize}
\item Height and weight
\item Rainfall and temperature
\item Vehicle speed and traffic density
\item Rolling two dice
\end{itemize}

When two quantities vary together, we need probability involving both variables.

Random variables may have a joint distribution.

If variables are independent, computation is simple.
If not independent, joint analysis is required.

Therefore, joint random variables are used to study the relation between variables.


\section{Discrete Case --- Joint PMF}

Suppose $X$ and $Y$ are discrete random variables.

\[
X = \{x_1, x_2, \dots , x_n\}
\]

\[
Y = \{y_1, y_2, \dots , y_m\}
\]

Then the ordered pair $(X,Y)$ takes values in

\[
\{(x_1,y_1), (x_1,y_2), \dots , (x_n,y_m)\}
\]

\subsection{Definition}

\[
p(x_i,y_j) = P(X=x_i , Y=y_j)
\]

This gives probability that both events occur simultaneously.

\subsection{Properties}

\[
0 \le p(x_i,y_j) \le 1
\]

\[
\sum_{i=1}^{n} \sum_{j=1}^{m} p(x_i,y_j) = 1
\]

All probabilities must be non-negative and total probability must be 1.


\section{Example --- Rolling Two Dice}

Roll two fair dice.

Let

\[
X = \text{value on first die}
\]

\[
T = \text{sum of both dice}
\]

Each outcome has probability

\[
\frac{1}{36}
\]

$X$ and $T$ are not independent because $T$ depends on $X$.

For example, if $X=6$, then $T \ge 6$.


\section{Continuous Case --- Joint PDF}

For continuous variables,

\[
X \in [a,b], \quad Y \in [c,d]
\]

Then

\[
(X,Y) \in [a,b] \times [c,d]
\]

\subsection{Definition}

A function $f(x,y)$ is joint density if

\[
P((X,Y) \text{ in small region})
\approx f(x,y) dx dy
\]

\subsection{Properties}

\[
f(x,y) \ge 0
\]

\[
\int_c^d \int_a^b f(x,y) dx dy = 1
\]

Density may be greater than 1 but total probability must be 1.


\section{Example}

Given

\[
f(x,y) = 4xy
\]

for

\[
0 \le x \le 1, \quad 0 \le y \le 1
\]

\subsection{Check Valid PDF}

\[
\int_0^1 \int_0^1 4xy dx dy
\]

\[
= \int_0^1 4y \left[\frac{x^2}{2}\right]_0^1 dy
\]

\[
= \int_0^1 2y dy
\]

\[
= [y^2]_0^1 = 1
\]

So valid PDF.


\subsection{Find Probability}

Event:

\[
A = \{X < 0.5 , Y > 0.5\}
\]

\[
P(A)
=
\int_0^{0.5}
\int_{0.5}^{1}
4xy \, dy \, dx
\]

\[
=
\int_0^{0.5}
2x(1-0.25) dx
\]

\[
=
\int_0^{0.5}
\frac{3}{2}x dx
\]

\[
=
\frac{3}{16}
\]


\section{Joint CDF}

Definition:

\[
F(x,y) = P(X \le x , Y \le y)
\]

CDF gives accumulated probability.


\section{Continuous Case}

If PDF is $f(x,y)$,

\[
F(x,y)
=
\int_c^y
\int_a^x
f(u,v) du dv
\]

PDF can be recovered by

\[
f(x,y)
=
\frac{\partial^2}{\partial x \partial y}
F(x,y)
\]


\section{Discrete Case}

If PMF is $p(x_i,y_j)$,

\[
F(x,y)
=
\sum_{x_i \le x}
\sum_{y_j \le y}
p(x_i,y_j)
\]

CDF is obtained by adding probabilities.


\section{Properties of Joint CDF}

\begin{itemize}
\item $F(x,y)$ is non-decreasing
\item $F(x,y)=0$ at lower limit
\item $F(x,y)=1$ at upper limit
\end{itemize}


\section{Exam Recall}

Must remember:

\begin{itemize}
\item Joint PMF definition
\item PMF properties
\item Dice example
\item Joint PDF definition
\item Double integral = 1
\item Example with $4xy$
\item Joint CDF definition
\item Continuous CDF integral
\item Discrete CDF sum
\item PDF from CDF derivative
\item CDF properties
\end{itemize}

\end{document}